\section*{第 \ref{ch:g9:inside-the-atom} 章——原子内部}

\begin{solution}{exo:g9:inside-the-atom:1}
\omterm{def:g9:inside-the-atom:nucleus}{原子核}中有\omterm{def:g9:inside-the-atom:proton}{质子}（带正电）和\omterm{def:g9:inside-the-atom:proton}{中子}（不带电），
周围是\omterm{def:g9:inside-the-atom:nucleus}{电子}（带负电）。
\end{solution}

\begin{solution}{exo:g9:inside-the-atom:2}
\ce{^{16}_{8}O}：8 个\omterm{def:g9:inside-the-atom:proton}{质子}、8 个\omterm{def:g9:inside-the-atom:proton}{中子}、8 个\omterm{def:g9:inside-the-atom:nucleus}{电子}。
\ce{^{27}_{13}Al}: 13, 14, 13. \ce{^{40}_{20}Ca}: 20, 20, 20.
\end{solution}

\begin{solution}{exo:g9:inside-the-atom:3}
它的\omterm{def:g9:inside-the-atom:nucleus}{电子}（每个带电荷 $-e$）与\omterm{def:g9:inside-the-atom:proton}{质子}（每个带电荷 $+e$）
数目相等，电荷相互抵消。
\end{solution}

\begin{solution}{exo:g9:inside-the-atom:4}
$A = 26 + 30 = 56$: \ce{^{56}_{26}Fe}.
\end{solution}

\begin{solution}{exo:g9:inside-the-atom:5}
是的：两者的 $Z = 17$，因此都是氯。它们的
\omterm{def:g9:inside-the-atom:atomic-number}{质量数}不同（\omterm{def:g9:inside-the-atom:proton}{中子}数分别为 18 和 20），它们互为\omterm{def:g9:inside-the-atom:isotope}{同位素}。
\end{solution}

\begin{solution}{exo:g9:inside-the-atom:6}
否：它们的\omterm{def:g9:inside-the-atom:atomic-number}{原子序数}不同（分别为 6 和 7），因此是不同\omterm{def:g9:inside-the-atom:element}{元素}，
分别为碳和氮。\omterm{def:g9:inside-the-atom:isotope}{同位素}的 $Z$ 相同。
\end{solution}

\begin{solution}{exo:g9:inside-the-atom:7}
$56 \times 1.67 \times 10^{-27} \approx \qty{9.35e-26}{kg}$.
\end{solution}

\begin{solution}{exo:g9:inside-the-atom:8}
\omterm{def:g7:atoms-and-molecules:atom}{原子}主要是空的：\omterm{def:g9:inside-the-atom:nucleus}{原子核}很小，
大多数粒子在途中遇不到\omterm{def:g9:inside-the-atom:nucleus}{原子核}，直接穿过。
\end{solution}

\begin{solution}{exo:g9:inside-the-atom:9}
$10^{-10} / 10^{-15} = 10^{5}$：小了十万倍。
\end{solution}

\begin{solution}{exo:g9:inside-the-atom:10}
化学性质取决于\omterm{def:g9:inside-the-atom:nucleus}{电子}，而同一\omterm{def:g9:inside-the-atom:element}{元素}的\omterm{def:g9:inside-the-atom:isotope}{同位素}
\omterm{def:g9:inside-the-atom:nucleus}{电子}数相同（即 $Z$ 相同）。
\end{solution}

\begin{solution}{exo:g9:inside-the-atom:11}
$26 \times 9.11 \times 10^{-31} = \qty{2.37e-29}{kg}$。占比：
$2.37 \times 10^{-29} / 9.35 \times 10^{-26} \approx 2.5 \times
10^{-4}$，约占质量的 \qty{0.025}{\percent}。
\end{solution}

\begin{solution}{exo:g9:inside-the-atom:12}
6915 个铜-63 \omterm{def:g7:atoms-and-molecules:atom}{原子}和 3085 个铜-65 \omterm{def:g7:atoms-and-molecules:atom}{原子}。平均值：
$(6915 \times 63 + 3085 \times 65)/10\,000 = 63.6$ 个\omterm{def:g9:inside-the-atom:proton}{核子}。
\end{solution}

\begin{solution}{pb:g9:inside-the-atom:1}
\textbf{1.} \ce{^{35}_{17}Cl} 和 \ce{^{37}_{17}Cl}。

\textbf{2.} 氯-35：17 个\omterm{def:g9:inside-the-atom:proton}{质子}、18 个\omterm{def:g9:inside-the-atom:proton}{中子}、17 个\omterm{def:g9:inside-the-atom:nucleus}{电子}。
氯-37：17 个\omterm{def:g9:inside-the-atom:proton}{质子}、20 个\omterm{def:g9:inside-the-atom:proton}{中子}、17 个\omterm{def:g9:inside-the-atom:nucleus}{电子}。

\textbf{3.} 两者都有 17 个\omterm{def:g9:inside-the-atom:proton}{质子}，因此都是氯\omterm{def:g9:inside-the-atom:element}{元素}；它们
仅\omterm{def:g9:inside-the-atom:proton}{中子}数不同，因此互为\omterm{def:g9:inside-the-atom:isotope}{同位素}。

\textbf{4.} 是的：化学性质取决于\omterm{def:g9:inside-the-atom:nucleus}{电子}，而两者的\omterm{def:g9:inside-the-atom:nucleus}{电子}数都是 17。

\textbf{5.} $35 \times 1.67 \times 10^{-27} = \qty{5.845e-26}{kg}$.

\textbf{6.} $17 \times 9.11 \times 10^{-31} = \qty{1.55e-29}{kg}$:
比\omterm{def:g9:inside-the-atom:nucleus}{原子核}约小 4000 倍，可忽略不计。

\textbf{7.} $37 \times 1.67 \times 10^{-27} = \qty{6.179e-26}{kg}$.

\textbf{8.} $\qty{1}{g} = \qty{e-3}{kg}$; $10^{-3} / 5.845 \times
10^{-26} \approx 1.71 \times 10^{22}$ 个\omterm{def:g7:atoms-and-molecules:atom}{原子}。

\textbf{9.} \qty{75.8}{\percent} 和 \qty{24.2}{\percent}。

\textbf{10.} $758 \times 5.845 \times 10^{-26} = \qty{4.430e-23}{kg}$;
$242 \times 6.179 \times 10^{-26} = \qty{1.495e-23}{kg}$.

\textbf{11.} $4.430 \times 10^{-23} + 1.495 \times 10^{-23} =
\qty{5.925e-23}{kg}$.

\textbf{12.} $5.925 \times 10^{-23} / 1000 \approx
\qty{5.93e-26}{kg}$：一个氯\omterm{def:g7:atoms-and-molecules:atom}{原子}的平均质量。
\end{solution}
